\documentclass[10pt]{article}
\usepackage{rikai-textbook}
\begin{document}
\rikaititle{sample-doc.md}{2026-10-04}
\tableofcontents

\section{2. 受付}
% >>> topic-001
\topic{一意キーと重複の防止}{復習中}
同じ送信要求が 2 度届いたとき、ジョブが増えるかどうかを判断する論点である。

\begin{youten}
送信要求には、呼び出し元が利用者 ID と通知種別から作った一意キーを付ける。同じ一意キーを持つ要求を再度受け付けた場合、ワーカーは新しいジョブを作らず、既存ジョブの ID を返す。

この規則が成り立つ前提は、呼び出し元が最初の要求と再送要求に同じ一意キーを使うことである。キーを作り直した要求は、別のジョブとして扱う。
\end{youten}

\begin{tabular}{L{0.36}L{0.54}}
\toprule
再送要求の一意キー & ワーカーの扱い \\
\midrule
最初の要求と同じキー & 新しいジョブを作らず、既存ジョブの ID を返す \\
作り直したキー & 別のジョブとして扱う \\
\bottomrule
\end{tabular}

\begin{reidai}[同じキーで再送する場合（説明用）]
呼び出し元が一意キー $K$ を付けた要求を送り、ワーカーがジョブ $J$ として受け付けた。その後、同じ要求が届いた。
\begin{enumerate}
  \item 再送要求のキーが $K$ のままなら、受け付け済みと同じキーである。ワーカーは新しいジョブを作らず、$J$ の ID を返す。
  \item 再送要求のキーを $K'$ に作り直していたら、ワーカーは $J$ とは別のジョブとして扱う。
\end{enumerate}
\end{reidai}

\begin{mayoi}
「新しいジョブを作る」のは、キーを作り直した要求のときだけである。同じキーのままなら、何度受け付けてもジョブは増えず、返る ID は既存ジョブのものである。見分ける条件は、受付時のキーが同じかどうかである。
\end{mayoi}

\begin{genbun}
「同じキーを持つ要求を再度受け付けた場合、ワーカーは新しいジョブを作らず、既存ジョブの ID を返す。」\\
「キーを作り直した要求は別のジョブとして扱う。」
\end{genbun}
% <<< topic-001

% >>> topic-002
\topic{配信順序は保証されない}{卒業}
同じ利用者へのジョブについて、保証される順序と保証されない順序を区別する論点である。

\begin{youten}
ワーカーは同じ利用者へのジョブを受付順に保存する。配信の完了順は保証しない。先に受け付けたジョブが再送待ちになると、後のジョブが先に完了することがある。
\end{youten}

\begin{tabular}{lL{0.6}}
\toprule
順序 & 原文の定め \\
\midrule
保存の順 & 受付順に保存する \\
完了の順 & 保証しない。再送待ちが入ると入れ替わりうる \\
\bottomrule
\end{tabular}

\begin{center}
\begin{tikzpicture}[job/.style={draw,circle,minimum size=8mm},>=Stealth]
  \node (l1) {保存の順};
  \node[job,right=8mm of l1] (a1) {A};
  \node[job,right=6mm of a1] (b1) {B};
  \node[below=12mm of l1] (l2) {完了の順};
  \node[job,right=20mm of l2] (b2) {B};
  \node[job,right=24mm of b2] (a2) {A};
  \draw[->] (b1) -- (b2);
  \draw[->,dashed] (a1) -- node[above right,font=\small] {再送待ち} (a2);
\end{tikzpicture}
\end{center}

\begin{reidai}[先のジョブが再送待ちになる場合（説明用）]
同じ利用者へのジョブ A、B を、この順に受け付けた。A だけが再送待ちになった。
\begin{enumerate}
  \item 保存の順は受付順なので、A、B の順である。再送待ちになっても、この順は変わらない。
  \item A は再送待ちの間は完了しない。そのため、B が A より先に完了することがある。
  \item 完了の順は保証されないので、B が先に完了しても原文の規則に反しない。
\end{enumerate}
\end{reidai}

\begin{mayoi}
受付順が守られるのは保存までで、完了の順には及ばない。「受付順に保存した」ことは、後のジョブが先に完了しない理由にはならない。後のジョブが先に完了しうるのは、先のジョブが再送待ちになった場合である。
\end{mayoi}

\begin{genbun}
「ワーカーは同じ利用者へのジョブを受付順に保存するが、配信の完了順は保証しない。先に受け付けたジョブが再送待ちになると、後のジョブが先に完了することがある。」
\end{genbun}
% <<< topic-002

\section{3. 再送}
% >>> topic-003
\topic{再送対象の判定（502/503/タイムアウト）}{復習中}
外部 API への送信が失敗したとき、失敗の種類によって再送するかどうかを決める論点である。

\begin{youten}
外部 API が \texttt{502} または \texttt{503} を返した場合と、接続がタイムアウトした場合は再送する。これらは一時的な失敗として扱う。

一方、\texttt{429} を除く \texttt{4xx} 応答は再送せず、失敗として確定する。要求内容を変えないまま再送しても成功が見込めないためである。
\end{youten}

\begin{tabular}{L{0.34}ll}
\toprule
失敗の種類 & 扱い & 処理 \\
\midrule
\texttt{502} 応答 & 一時的な失敗 & 再送する \\
\texttt{503} 応答 & 一時的な失敗 & 再送する \\
接続のタイムアウト & 一時的な失敗 & 再送する \\
\texttt{4xx} 応答（\texttt{429} を除く） & 失敗として確定 & 再送しない \\
\bottomrule
\end{tabular}

\begin{reidai}[3 つの失敗を仕分ける（説明用）]
ジョブ X には \texttt{503} が、ジョブ Y には接続のタイムアウトが、ジョブ Z には \texttt{404} が返った。
\begin{enumerate}
  \item X：\texttt{503} は一時的な失敗なので、再送する。
  \item Y：接続のタイムアウトは一時的な失敗なので、再送する。
  \item Z：\texttt{404} は \texttt{429} 以外の \texttt{4xx} なので、再送せずに失敗として確定する。
\end{enumerate}
\end{reidai}

\begin{mayoi}
応答が返ったかどうかでは分かれない。\texttt{503} は応答が返っていても、一時的な失敗として再送する。また、エラー応答がすべて再送の対象になるわけではなく、\texttt{429} を除く \texttt{4xx} は再送しない。分かれ目は応答コードである。
\end{mayoi}

\begin{genbun}
「外部 API が \texttt{502} または \texttt{503} を返した場合と、接続がタイムアウトした場合は再送する。これらは一時的な失敗として扱う。」
\end{genbun}
% <<< topic-003

% >>> topic-004
\topic{429 は「再送しない」の例外}{卒業}
\texttt{4xx} 応答を再送しないという通常則と、その例外を区別する論点である。

\begin{youten}
\texttt{429} を除く \texttt{4xx} 応答は再送しない。要求内容を変えないまま再送しても成功が見込めないため、失敗として確定する。\texttt{429} は、この「再送しない」の規則から除かれる。原文は、\texttt{429} を再送するかどうかを明記していない。
\end{youten}

\begin{tabular}{lL{0.56}}
\toprule
応答 & 「\texttt{4xx} は再送しない」の規則 \\
\midrule
\texttt{429} 以外の \texttt{4xx} & 当てはまる。再送せず、失敗として確定する \\
\texttt{429} & 当てはまらない（例外） \\
\bottomrule
\end{tabular}

\begin{genbun}
「\texttt{429} を除く \texttt{4xx} 応答は再送しない。要求内容を変えないまま再送しても成功が見込めないため、失敗として確定する。」
\end{genbun}
% <<< topic-004

% >>> topic-005
\topic{試行上限と失敗キュー}{卒業}
再送の前にどれだけ待ち、何回で自動の再送を打ち切るかを定める論点である。

\begin{youten}
再送前の待機時間は 1 秒、2 秒、4 秒の順に延ばし、その後は 8 秒を上限とする。試行回数は、初回送信を含めて最大 5 回である。5 回目が失敗したジョブは自動では再送せず、失敗キューへ移す。

初回送信も試行に数えるので、再送の回数は最大で次のとおりである。
\begin{equation}
  \text{再送の回数} = \text{試行回数の上限} - \text{初回送信} = 5 - 1 = 4
\end{equation}
\end{youten}

\begin{tabular}{lL{0.56}}
\toprule
試行 & 直前の待機時間 \\
\midrule
1 回目（初回送信） & なし（再送ではない） \\
2 回目 & 1 秒 \\
3 回目 & 2 秒 \\
4 回目 & 4 秒 \\
5 回目 & 上限の 8 秒以内 \\
5 回目も失敗 & 自動では再送せず、失敗キューへ移す \\
\bottomrule
\end{tabular}

\begin{reidai}[5 回とも失敗したときの待機時間の合計]
待機は再送の前にだけ入るので、4 回ある。4 回目の待機は上限の 8 秒以内なので、待機時間の合計 $\mathrm{T}$ は最大で次のとおりである。送信そのものにかかる時間は含めない。
\begin{align}
  \mathrm{T} &\leq 1 + 2 + 4 + 8 \\
             &= 15 \text{ 秒}
\end{align}
その後、ジョブは失敗キューへ移り、自動では再送されない。
\end{reidai}

\begin{genbun}
「再送前の待機時間は 1 秒、2 秒、4 秒の順に延ばし、その後は 8 秒を上限とする。」\\
「初回送信を含む試行回数は最大 5 回とする。5 回目が失敗したジョブは、自動では再送せず、失敗キューへ移す。」
\end{genbun}
% <<< topic-005

\section{4. 完了と停止}
% >>> topic-006
\topic{完了判定と保存失敗}{復習中}
ジョブを完了とする条件を判断する論点である。条件は 2 つあり、両方がそろって初めて完了になる。

\begin{youten}
外部 API が成功応答を返し、その応答 ID を配信記録へ保存できた時点で、ジョブを完了とする。成功応答を受け取っても保存に失敗した場合は、完了にしない。
\begin{equation}
  \text{完了} \iff \text{成功応答を受け取った} \land \text{応答 ID を保存できた}
\end{equation}
保存に失敗したジョブを再実行するときも、受付時の一意キーを引き継ぐ。これにより、外部 API への重複送信の有無を追跡できる。
\end{youten}

\begin{tabular}{lll}
\toprule
成功応答 & 応答 ID の保存 & ジョブ \\
\midrule
受け取った & 成功 & 完了 \\
受け取った & 失敗 & 完了にしない \\
\bottomrule
\end{tabular}

\begin{reidai}[成功応答の後に保存へ失敗した場合]
外部 API から成功応答を受け取ったが、応答 ID を配信記録へ保存できなかった。
\begin{enumerate}
  \item 完了の条件のうち、保存ができていない。よって、ジョブは完了にしない。
  \item このジョブを再実行するときは、受付時の一意キーを引き継ぐ。
  \item キーが同じなので、外部 API への重複送信の有無を追跡できる。
\end{enumerate}
\end{reidai}

\begin{mayoi}
外部 API の成功応答は、完了の条件の片方にすぎない。もう片方は、応答 ID を配信記録へ保存できたことである。成功応答を受け取っていても、保存に失敗すれば完了ではない。
\end{mayoi}

\begin{genbun}
「外部 API が成功応答を返し、その応答 ID を配信記録へ保存できた時点で、ジョブを完了とする。成功応答を受け取っても保存に失敗した場合は完了にしない。」
\end{genbun}
% <<< topic-006

% >>> topic-007
\topic{緊急停止の例外}{卒業}
運用者が配信を止めたとき、すでに送信中の要求をどう扱うかを、通常停止と緊急停止で区別する論点である。

\begin{youten}
運用者が配信を停止すると、ワーカーは新しいジョブの取得を止める。すでに外部 API へ送信中の要求は取り消さず、その応答を記録してからワーカーを停止する。

ただし、情報漏えいのおそれがある緊急停止では、送信中の要求も直ちに中断する。通常停止の完了待ちは適用しない。
\end{youten}

\begin{tabular}{L{0.34}lL{0.32}}
\toprule
停止の種類 & 送信中の要求 & ワーカーの停止 \\
\midrule
通常停止 & 取り消さない & 応答を記録してから停止する \\
緊急停止（情報漏えいのおそれがある） & 直ちに中断する & 完了待ちを適用しない \\
\bottomrule
\end{tabular}

\begin{genbun}
「すでに外部 API へ送信中の要求は取り消さず、その応答を記録してからワーカーを停止する。」\\
「ただし、情報漏えいのおそれがある緊急停止では、送信中の要求も直ちに中断する。通常停止の完了待ちは適用しない。」
\end{genbun}
% <<< topic-007

\end{document}
